\section{Wind Farm Type I: Kusiak--Song Benchmark}
\label{sec:results}
\label{sec:overall}
\label{sec:farm-I}
Type~I is the Kusiak--Song benchmark~\cite{Kusiak2010}: 68 circular farms under two synthetic wind roses, on which the nine methods are compared at one budget. This section gives the experimental setting, the numerical results and their statistical analysis (including the primary comparison of PSO-VNS with PSO and an all-run endpoint), followed by studies of the budget and initialization (Section~\ref{sec:beyond}) and of the model assumptions (Section~\ref{sec:robustness}).

\subsection{Experimental Setting}
The setting is that of Type~I in Table~\ref{tab:farm-types}:
\begin{enumerate}
\item \emph{Cases}: 68, the two wind data sets combined with three farm radii, $r=500$~m with $N=2,\ldots,10$ turbines (9 cases per data set), $r=750$~m with $N=2,\ldots,12$ (11) and $r=1000$~m with $N=2,\ldots,15$ (14).
\item \emph{Wind data} (Section~\ref{sec:wind}; Table~\ref{S-tab:S-wind}): Data Set~I, a nearly unidirectional rose with direction-independent wind speeds (Weibull scale 13~m/s in all directions), and Data Set~II, a broad rose with direction-dependent wind speeds (scales 2.6--10~m/s); both with 24 direction bins of 15$^\circ$, Weibull shape 2 and 21 speed bins of 0.5~m/s between cut-in and rated speed.
\item \emph{Model} (Section~\ref{sec:model}): the Jensen top-hat wake ($K=0.075$, $C_T=0.8$, hub-center test, root-sum-square superposition) and the linear Kusiak--Song power curve (cut-in 3.5~m/s, rated 14~m/s, 1500~kW; $D=77$~m); the objective is the wake loss in percent of the wake-free value.
\item \emph{Constraints} (Section~\ref{sec:constraints}): minimum spacing $4D=308$~m and the circular boundary of radius $r$, handled by the penalty~\eqref{eq:penalty} with box clipping to $[-r,r]^{2N}$.
\item \emph{Runs}: 6{,}030 objective evaluations per run, random initialization and 30 seed-paired runs (seeds 1--30) per case and method, 18{,}360 runs in total.
\item \emph{Methods}: the nine methods of Section~\ref{sec:setup-methods} with the settings of Table~\ref{tab:params}.
\item \emph{Analysis}: the feasibility-aware ranking and the tests of Section~\ref{sec:setup-stats}; per-case results in Tables~\ref{S-tab:res-1-500}--\ref{S-tab:res-2-1000} (GA: Table~\ref{S-tab:S-xp-ga}).
\end{enumerate}

\subsection{Numerical Results}
% generated by analysis/rev4_farm_tables.py; regenerate, do not edit
\begin{table}[!htb]
\centering
\caption{Type~I farms (Kusiak--Song benchmark): mean wake loss (\%) by data set and farm radius, feasible runs and average rank of the nine methods.}
\label{tab:farm-I-res}
\scriptsize\setlength{\tabcolsep}{2.6pt}
\begin{tabular}{lccccccccc}
\toprule
& \multicolumn{3}{c}{Data Set I} & \multicolumn{3}{c}{Data Set II} & & & \\
\cmidrule(lr){2-4}\cmidrule(lr){5-7}
Method & 500~m (9) & 750~m (11) & 1000~m (14) & 500~m (9) & 750~m (11) & 1000~m (14) & All (68) & Feas.\ (\%) & Avg.\ rank \\
\midrule
PSO-VNS & 1.035 & 0.522 & 0.463 & \textbf{2.923} & \textbf{1.856} & \textbf{1.547} & \textbf{1.322} & \textbf{100.0} & \textbf{2.18} \\
PSO & \textbf{0.971} & 0.479 & 0.471 & 2.986 & 1.887 & 1.637 & 1.341 & \textbf{100.0} & 2.40 \\
SSA-VNS & 1.503 & 0.771 & 0.619 & 3.364 & 2.235 & 1.919 & 1.653 & 99.6 & 4.32 \\
SSA & 1.363$^{8}$ & 0.991 & 0.833 & 3.101$^{8}$ & 2.575 & 2.324 & 1.805$^{66}$ & 97.6 & 6.34 \\
LX-SSA & 1.607$^{8}$ & 1.116 & 0.941 & 3.330$^{8}$ & 2.746 & 2.437 & 1.959$^{66}$ & 96.9 & 7.01 \\
DE & 0.690$^{6}$ & 0.855$^{8}$ & 0.977$^{11}$ & 2.115$^{6}$ & 1.981$^{8}$ & 2.215$^{11}$ & 1.493$^{50}$ & 71.4 & 7.96 \\
GA & 1.062 & \textbf{0.476} & \textbf{0.441} & 3.071 & 1.864 & 1.612 & 1.348 & 98.6 & 2.74 \\
VNS & 1.548 & 0.902 & 0.715 & 3.450 & 2.424 & 2.142 & 1.788 & \textbf{100.0} & 5.21 \\
MS-SLSQP & 1.180 & 0.887 & 0.854 & 3.355 & 2.360 & 2.219 & 1.758 & \textbf{100.0} & 6.85 \\
\bottomrule
\end{tabular}
\par\vspace{2pt}\parbox{\columnwidth}{\scriptsize 68 benchmark cases (number of cases per cluster in parentheses), 30 seed-paired runs per case and method, 6{,}030 evaluations, random initialization. Wake loss: per case the mean wake loss of the feasible runs, averaged over the cases of the cluster in which the method has at least 15 of 30 feasible runs (superscript: number of such cases when fewer than all; such means are over easier case subsets and are not comparable). Bold: lowest mean among the methods qualified in every case of the column; highest share of feasible runs; best average rank. Avg.\ rank: feasibility-aware rank of the mean feasible objective within the nine-method pool (ranking rule of Section~\ref{sec:setup-stats}), averaged over the 68 cases.}
\end{table}


Table~\ref{tab:farm-I-res} reports the mean wake loss of the nine methods for each farm size (500, 750 and 1000~m) under both wind data sets and over all 68 cases, with the share of feasible runs and the average rank. PSO-VNS has the best average rank, followed by PSO and GA, and the lowest mean wake loss in the three Data Set~II clusters and over all cases; PSO has the lowest in Data Set~I at 500~m and GA at 750 and 1000~m. PSO-VNS ends feasible in all runs (DE: 71.4\%); in the densest cases (500~m, $N=10$), the global best of its PSO phase is feasible at the switch in only 63.3\% of the runs, and the VNS phase makes all runs feasible. The wake loss grows with $N$, falls as the farm becomes larger, and is higher for the broader wind rose of Data Set~II.
% CHECK-FINAL [C42]: case-mean wake loss grows with N, falls with r, DS II >= DS I; [C08]: feasible_pct.PSOBV == 100; [C24]: dense_500_10 feasible at switch < 100

\subsection{Statistical Analysis}
% generated by analysis/rev4_farm_tables.py; regenerate, do not edit
\begin{table}[!htb]
\centering
\caption{Type~I farms (Kusiak--Song benchmark): Friedman test of the nine methods and comparisons of PSO-VNS with each method.}
\label{tab:farm-I-stats}
\footnotesize\setlength{\tabcolsep}{4pt}
\begin{tabular}{lcccccc}
\toprule
Method & Avg.\ rank & $p_z$ & Case means & $p_W$ & $\overline{\Delta L}$ (pp) & W/T/L \\
\midrule
PSO-VNS & 2.18 & -- & -- & -- & -- & -- \\
PSO & 2.40 & 0.639 & 34/8/26 & 0.296 & $-0.018$ & 13/52/3 \\
SSA-VNS & 4.32 & $1.6\times10^{-5}$ & 59/8/1 & $7.6\times10^{-11}$ & $-0.330$ & 49/19/0 \\
SSA & 6.34 & $4.0\times10^{-18}$ & 62/6/0 & $5.3\times10^{-11}$ & $-0.611$ & 59/9/0 \\
LX-SSA & 7.01 & $4.9\times10^{-24}$ & 62/6/0 & $5.3\times10^{-11}$ & $-0.765$ & 60/8/0 \\
DE & 7.96 & $6.8\times10^{-34}$ & 60/8/0 & $7.0\times10^{-11}$ & $-0.887$ & 56/12/0 \\
GA & 2.74 & 0.468 & 39/9/20 & 0.101 & $-0.026$ & 9/58/1 \\
VNS & 5.21 & $4.0\times10^{-10}$ & 55/9/4 & $1.9\times10^{-10}$ & $-0.465$ & 51/16/1 \\
MS-SLSQP & 6.85 & $1.4\times10^{-22}$ & 66/0/2 & $4.8\times10^{-11}$ & $-0.435$ & 58/10/0 \\
\bottomrule
\end{tabular}
\par\vspace{2pt}\parbox{\columnwidth}{\scriptsize 68 cases, 30 seed-paired runs, 6{,}030 evaluations. Friedman test on the feasibility-aware case ranks of the nine methods: $\chi^2_F=374.8$ (8 d.f.), $p=4.5\times10^{-76}$; Iman--Davenport $F_F=148.5$ (8, 536 d.f.), $p=1.2\times10^{-130}$. $p_z$: Holm-adjusted $p$ of the average-rank $z$ test against PSO-VNS (8 comparisons). Case means: cases in which PSO-VNS has the lower / equal (or neither method qualifies) / higher mean wake loss (a case in which only one method has at least 15 feasible runs counts for that method). $p_W$: two-sided Wilcoxon signed-rank test on the 68 per-case mean wake losses, PSO-VNS vs.\ method (cases in which only one method qualifies imputed in its favour), Holm-adjusted over the 8 comparisons. $\overline{\Delta L}$: mean difference of the case-mean wake losses, PSO-VNS minus method (percentage points; negative = PSO-VNS better), over the cases in which both qualify (66 for SSA, 66 for LX-SSA, 50 for DE, otherwise 68). W/T/L: cases in which PSO-VNS is significantly better / not different / worse in the seed-paired Wilcoxon test of the 30 runs (infeasible runs below every feasible run; Holm over the 8 comparisons of PSO-VNS in each case, $\alpha=0.05$).}
\end{table}
The Friedman test rejects equal performance of the nine methods (Table~\ref{tab:farm-I-stats}). PSO-VNS is significantly better than the other six methods in both the average-rank test and the case-mean Wilcoxon test, but not significantly different from PSO or GA. GA and PSO are practically equivalent at $\pm0.05$~pp at the case level (exploratory), whereas equivalence of GA and PSO-VNS is not shown (90\% intervals of GA minus PSO and GA minus PSO-VNS: [$-$0.022, 0.041] and [$-$0.006, 0.062]~pp; Table~\ref{S-tab:S-xp-ga}).

\subsubsection{PSO-VNS versus PSO}
\label{sec:psovns-pso}
Whether the VNS phase adds a practically relevant gain to PSO is a primary question (Section~\ref{sec:setup-stats}). It is decided by equivalence at the margin $m=\pm0.05$~pp (Definition~\ref{def:S-equiv}; Table~\ref{tab:equiv-main} and Fig.~\ref{fig:equiv-curve}); Table~\ref{tab:psovns-pso} gives the further tests and the subsets of cases. At the run level PSO-VNS has more significant wins than losses (13/52/3; Table~\ref{tab:farm-I-stats}), but the case means do not differ significantly. The findings are:
\begin{enumerate}
\item \emph{Equivalent on average.} The mean difference, $-0.018$~pp, has a 90\% interval inside $\pm0.05$~pp at the seed level, at which the verdict is decided, and at the case level, but not at the cluster level (six scenario clusters). The seed-level interval excludes zero: the two methods differ reproducibly on this benchmark, but by less than the margin (0.020\% of AEP). Practical equivalence is also the most probable outcome of the Bayesian signed-rank test.
\item \emph{Differences in both directions.} 19 cases differ by more than the margin, 11 in favor of PSO-VNS and 8 in favor of PSO; the two methods are equivalent for $N<10$, but not for $N\ge10$. Post hoc, PSO-VNS has the lower mean wake loss in all 10 Data Set~II cases with $N\ge10$, whereas the 7 Data Set~I cases beyond the margin, dense layouts in the two smaller farms, all favor PSO; this crossover has no identified cause and is a hypothesis for designed experiments.
\item \emph{Second phase.} After the switch both methods keep improving (Fig.~\ref{S-fig:S-conv-max}): the VNS phase removes slightly more of the remaining wake loss than continuing PSO for the same 3,000 evaluations, and it repairs the 26 runs infeasible at the switch.
\end{enumerate}
PSO-VNS thus does not dominate the standard PSO: it is at least as good on average, gains in larger Data Set~II layouts and loses in dense Data Set~I layouts. With 1$^\circ$ direction bins, its advantage is small but significant and equivalence is not shown (Table~\ref{tab:equiv-main}); optimized directly with 1$^\circ$ bins, its advantage exceeds the margin (Section~\ref{sec:direction}).
% CHECK-FINAL [C04]: PSOC: W >= L, mean loss PSOBV < PSOC, case-mean p >= 0.05; [C49]: case-level 90% CI inside +-m, P(rope most probable) printed via \NBayPSOVNSvsPSORope; [C52]: SD p and worst-run p >= 0.05; [C10]: phase2 PSOBV.mean >= PSOC_continued.mean > 0 ("both keep improving")
% CHECK-EXTRA [X39] seed and case level equivalent, cluster level not; [X40] seed-level CI excludes 0 inside m; [X41] minimal margins seed < case < m < cluster; [X48] Bayesian wording + posterior means; [X50] HL near 0; [X56] margin 4.1/2.1 MWh/yr per turbine, mean gain 0.2; [X25] energy < 0.05 % AEP
% CHECK-EXTRA [X43] 11/8 beyond +-m, max 0.41/0.27; [X44] N < 10 equivalent, N >= 10 not; [X45] non-trivial cases not equivalent; [X47] data set x density crossover (exploratory, post hoc); [X19] DS II N >= 10 survives all-family Holm ("<=", R5 O13); [X27] DS II N >= 10 ~0.2 % AEP
% CHECK-FINAL [C32]: DS II N >= 10 case-mean wins > losses
% CHECK-DIR [F10]: 1-deg PSO-VNS - PSO CI below 0, not equivalent; [F14]: DS II N >= 10 advantage survives at 1 deg

\begin{table}[!htb]
\centering
\caption{PSO-VNS versus PSO on the benchmark: further tests (top) and subsets of cases (bottom). $\Delta L$: case-mean wake loss of PSO-VNS minus that of PSO (pp; negative = PSO-VNS better); margin $m=0.05$~pp.}
\label{tab:psovns-pso}
\footnotesize\setlength{\tabcolsep}{3pt}
\begin{tabular}{@{}>{\raggedright\arraybackslash}p{0.585\textwidth}>{\raggedright\arraybackslash}p{0.40\textwidth}@{}}
\toprule
\multicolumn{2}{@{}l}{\emph{Further tests (68 cases)}} \\
Case-mean Wilcoxon test $p$; Hodges--Lehmann estimate & 0.30; $-0.003$~pp \\
Per-case SD of the runs; worst run (Wilcoxon $p$) & 0.55; 0.80 \\
Case-level TOST $p$; shifted-null check on this dataset & 0.011; 0.015 \\
Minimal margin $m_{\min}$, seed / case / cluster level & 0.026 / 0.041 / 0.067~pp \\
Bayesian signed-rank test, posterior means of the shares: PSO-VNS better by more than $m$ / within $\pm m$ / PSO better by more than $m$ & 0.25 / 0.62 / 0.12 ($P_{\rm rope}>0.99$) \\
Margin $m$ in energy per turbine, Data Set~I / II & 4.1 / 2.1~MWh/yr \\
Mean gain of PSO-VNS per turbine & 0.2~MWh/yr (0.020\% of AEP) \\
Share of the wake loss left at the switch removed in the last 3,000 evaluations, mean (median) & VNS phase 30.3\% (22.6\%); continued PSO 29.2\% (20.6\%) \\
Runs infeasible at the switch repaired by the VNS phase & 26 \\
\end{tabular}
\par\smallskip
\begin{tabular*}{\textwidth}{@{\extracolsep{\fill}}lccccc@{}}
\midrule
\emph{Subsets of cases} & $n$ & Beyond $-m$\,/\,$+m$ & $\overline{\Delta L}$ & Case 90\% CI & Equivalent \\
\midrule
All cases & 68 & 11/8 & $-0.018$ & [$-$0.041, 0.004] & yes \\
$N<10$ & 48 & 2/4 & $+0.007$ & [$-$0.005, 0.021] & yes \\
$N\ge10$ & 20 & 9/4 & $-0.079$ & [$-$0.142, $-$0.017] & no \\
Wake loss of PSO-VNS $\ge0.2$~pp & 44 & 11/8 & $-0.026$ & [$-$0.060, 0.008] & no \\
Data Set~II, $N\ge10^{a}$ & 10 & \multicolumn{4}{l}{PSO-VNS lower in 10 of 10; $p_{\rm Holm}\le0.018$; gain 0.221\% of AEP} \\
Data set $\times$ density$^{a}$ & 34 & \multicolumn{4}{l}{interaction $p=0.01$; the 7 Data Set~I cases beyond $m$ favor PSO} \\
\bottomrule
\end{tabular*}
\par\vspace{2pt}\parbox{\textwidth}{\scriptsize Beyond: cases with $\Delta L<-m$ (PSO-VNS better by more than the margin; up to 0.41~pp) / $\Delta L>+m$ (PSO better; up to 0.27~pp); Equivalent: case-level 90\% interval inside $\pm m$ (verdicts at the three levels: Table~\ref{tab:equiv-main}). $^{a}$Post hoc, exploratory. Data Set~II, $N\ge10$: case-mean Wilcoxon test, Holm over all 28 case-mean tests. Density $\phi=N(\ell_{\min}/2r)^2$, the nominal share of the farm disc covered by the exclusion discs; interaction test over the 34 pairs of cases with identical $(r,N)$ in the two data sets. Details: Tables~\ref{S-tab:X-heterogeneity}, \ref{S-tab:X-beyond}, \ref{S-tab:X-bayes} and~\ref{S-tab:X-hl}.}
\end{table}

% --- tab:X-equiv-levels and tab:equivalence (all 18 pairs) moved to the supplementary material; compact main-text summary:
\begin{table}[!t]
\centering
\caption{Practical equivalence at the margin $m=0.05$~pp (a post hoc convention) for the primary questions and the sampling-control contrasts (68 cases, 6,030 evaluations, 15$^\circ$ rose unless stated). $\overline{\Delta L}$: mean case-mean wake-loss difference, first minus second method (pp; negative = first better). 90\% intervals at the three levels of Definition~\ref{def:S-equiv}: seed (fixed benchmark; the level at which verdicts are decided), case and cluster (sensitivity analyses; CR2, six scenario clusters, 5 d.f.). Verdict at the seed / case / cluster level: E, equivalent (interval inside $\pm m$); NS, a gain of the first method larger than $m$ is ruled out (non-superiority); G, first method better by more than $m$; --, inconclusive. $P_{\rm rope}$: share of Bayesian signed-rank posterior samples in which practical equivalence is the most probable outcome. All 18 pairs, TOST $p$ values and minimal margins: Tables~\ref{S-tab:X-equiv-levels} and~\ref{S-tab:equivalence}.}
\label{tab:equiv-main}
\scriptsize\setlength{\tabcolsep}{2pt}
\begin{tabular}{@{}lcccccc@{}}
\toprule
Pair (first vs.\ second) & $\overline{\Delta L}$ & Seed 90\% CI & Case 90\% CI & Cluster 90\% CI & Verdict & $P_{\rm rope}$ \\
\midrule
\multicolumn{7}{@{}l}{\emph{Primary}} \\
PSO-VNS vs.\ PSO & $-0.018$ & [$-$0.026, $-$0.010] & [$-$0.041, 0.004] & [$-$0.067, 0.031] & E / E / -- & $>$0.99 \\
\quad 1$^\circ$ re-evaluation$^{a}$ & $-0.043$ & n/r & [$-$0.062, $-$0.025] & n/r & n/r / -- / n/r & n/r \\
SSA-VNS vs.\ RSD-VNS & $+0.024$ & [0.002, 0.046] & [$-$0.022, 0.062] & [$-$0.051, 0.099] & NS / NS / -- & 0.77 \\
\midrule
\multicolumn{7}{@{}l}{\emph{Exploratory: sampling controls}} \\
SSA-VNS vs.\ RS-VNS & $-0.068$ & [$-$0.092, $-$0.044] & [$-$0.101, $-$0.039] & [$-$0.126, $-$0.011] & -- / -- / -- & 0.70 \\
RSD-VNS vs.\ RS-VNS & $-0.093$ & [$-$0.115, $-$0.070] & [$-$0.127, $-$0.055] & [$-$0.131, $-$0.054] & G / G / G & 0.05 \\
LX-SSA-VNS vs.\ RSD-VNS & $+0.087$ & [0.065, 0.110] & [0.047, 0.126] & [0.004, 0.170] & NS / NS / NS & 0.10 \\
PSO-VNS vs.\ RSD-VNS & $-0.306$ & [$-$0.325, $-$0.287] & [$-$0.380, $-$0.239] & [$-$0.442, $-$0.171] & G / G / G & $<$0.01 \\
\bottomrule
\end{tabular}
\par\vspace{2pt}\parbox{\linewidth}{\scriptsize $^{a}$Stored 15$^\circ$ layouts re-evaluated with 1$^\circ$ sub-bins, not re-optimized (Section~\ref{sec:direction}; Table~\ref{S-tab:F-equiv}); n/r: not reported here. SSA-VNS vs.\ RSD-VNS at the seed level is also inside $\pm m$ (not under joint seed resampling, Table~\ref{S-tab:S-sa-syncseed}); the primary question for this pair is one-sided.}
\end{table}
% CHECK-EXTRA [X55] no pair equivalent at the cluster level; main-comparison pairs other than PSO: whole 90% CI beyond -m at all three levels (tab:X-equiv-levels, m_min > m)

\begin{figure}[!t]
\centering
\pipefig{0.45\textwidth}{figures_mpce/equiv_curve.pdf}{5cm}
\caption{Equivalence curve of PSO-VNS vs.\ PSO: TOST $p$ value versus the equivalence margin $m$ (pp of wake loss) at the three levels of inference of Table~\ref{tab:equiv-main}: seed level (bootstrap of the seed-paired runs within every case), case level (bootstrap over the 68 cases) and cluster level (six farm clusters; CR2 cluster-robust $t$ with 5 d.f., solid, and restricted wild-cluster bootstrap, dashed). Dotted lines: $\alpha=0.05$ and $m=0.05$~pp. Dots: minimal margins $m_{\min}$; equivalence holds at every margin to the right of a dot. The bootstrap curves stop at their resolution limit $1/10{,}001$.}
\label{fig:equiv-curve}
\end{figure}
% CHECK-EXTRA [X57] equiv_curve.pdf generated in the same run; [X41] minimal margins seed < case < m < cluster

\subsubsection{All-Run Reliability and Coordinate Precision}
\label{sec:archive-reliability}
The feasibility-aware rank compares conditional means and can favor a method with a better mean among fewer feasible runs: on Lillgrund at 6,030 evaluations, PSO-VNS has 17/30 feasible runs and a conditional mean AEP of 113.074~GWh, MS-SLSQP 30/30 and 112.865~GWh. We therefore add an all-run paired endpoint, in which for each seed pair
\begin{enumerate}
\item a feasible run beats an infeasible one;
\item of two feasible runs, the better objective wins (ties within $10^{-9}$);
\item two infeasible runs tie.
\end{enumerate}
Over $n_s$ seed pairs with $n_W$ wins and $n_T$ ties, the score is $(n_W+n_T/2)/n_s$ (0.5: no difference). Table~\ref{tab:allrun} gives it for the nine methods. On the benchmark, PSO-VNS is better more often than not against every method, also with the 30 synchronized seeds as the unit of replication; this frequency does not measure the size of the differences, and against PSO the mean difference lies within the margin (Table~\ref{tab:equiv-main}). SSA-VNS loses slightly more often than it wins against RSD-VNS. On Lillgrund, PSO-VNS is better than PSO at both budgets and than GA at 6,030 evaluations (at 30,030 the GA comparison is not significant), but not better than MS-SLSQP at 6,030 evaluations.

\begin{table}[!htb]
\centering
\caption{All-run paired outcome: score of the first method (0.5: no difference; above 0.5: first method better).}
\label{tab:allrun}
\scriptsize\setlength{\tabcolsep}{3pt}
\begin{tabular}{@{}lcccccc@{}}
\toprule
& & & \multicolumn{2}{c}{Pooled seed pairs} & \multicolumn{2}{c}{Seeds as unit} \\
\cmidrule(lr){4-5}\cmidrule(l){6-7}
Pair (first vs.\ second) & W/T/L & Score & 95\% CI & $p_{\rm Holm}$ & 95\% CI & $p_{\rm Holm}$ \\
\midrule
\multicolumn{7}{@{}l}{\emph{Benchmark, 68 cases, 2,040 seed pairs}} \\
PSO-VNS vs.\ PSO & 950/348/742 & 0.551 & [0.516, 0.587] & $4.7\times10^{-7}$ & [0.531, 0.572] & $1.2\times10^{-5}$ \\
PSO-VNS vs.\ GA & 969/318/753 & 0.553 & [0.524, 0.582] & $4.3\times10^{-7}$ & [0.536, 0.571] & $5.5\times10^{-6}$ \\
PSO-VNS vs.\ SSA-VNS & 1336/321/383 & 0.734 & [0.700, 0.766] & $6.5\times10^{-123}$ & [0.715, 0.752] & $1.5\times10^{-8}$ \\
PSO-VNS vs.\ SSA & 1650/276/114 & 0.876 & [0.839, 0.911] & $<10^{-300}$ & [0.866, 0.886] & $1.5\times10^{-8}$ \\
PSO-VNS vs.\ LX-SSA & 1658/273/109 & 0.880 & [0.842, 0.914] & $<10^{-300}$ & [0.868, 0.891] & $1.5\times10^{-8}$ \\
PSO-VNS vs.\ DE & 1682/296/62 & 0.897 & [0.851, 0.938] & $<10^{-300}$ & [0.888, 0.905] & $1.5\times10^{-8}$ \\
PSO-VNS vs.\ VNS & 1406/338/296 & 0.772 & [0.729, 0.812] & $3.8\times10^{-172}$ & [0.751, 0.792] & $1.5\times10^{-8}$ \\
PSO-VNS vs.\ MS-SLSQP & 1649/218/173 & 0.862 & [0.823, 0.897] & $<10^{-300}$ & [0.850, 0.874] & $1.5\times10^{-8}$ \\
SSA-VNS vs.\ RSD-VNS & 812/318/910 & 0.476 & [0.452, 0.500] & 0.019 & [0.455, 0.498] & 0.047 \\
\midrule
\multicolumn{7}{@{}l}{\emph{Lillgrund, 30 seed pairs}} \\
PSO-VNS vs.\ PSO, 6,030 & 17/13/0 & 0.783 & -- & -- & [0.700, 0.867] & $1.98\times10^{-4}$ \\
PSO-VNS vs.\ PSO, 30,030 & 28/0/2 & 0.933 & -- & -- & [0.833, 1.000] & $1.39\times10^{-5}$ \\
PSO-VNS vs.\ GA, 6,030 & 17/13/0 & 0.783 & -- & -- & [0.700, 0.867] & $1.98\times10^{-4}$ \\
PSO-VNS vs.\ GA, 30,030 & 21/0/9 & 0.700 & -- & -- & [0.533, 0.867] & 0.214 \\
PSO-VNS vs.\ MS-SLSQP, 6,030 & 10/0/20 & 0.333 & -- & -- & [0.167, 0.500] & 0.395 \\
PSO-VNS vs.\ MS-SLSQP, 30,030 & 28/0/2 & 0.933 & -- & -- & [0.833, 1.000] & $1.39\times10^{-5}$ \\
\bottomrule
\end{tabular}
\par\vspace{2pt}\parbox{\columnwidth}{\scriptsize Pooled seed pairs: 95\% interval from resampling the cases, exact sign test over the discordant pairs treated as independent. Seeds as unit: 95\% interval from resampling the 30 seeds (benchmark: complete seed vectors over all cases; 25 of 29 non-tied seeds favor PSO-VNS against PSO) and exact sign-flip (benchmark) or sign test (Lillgrund). Holm over the comparisons of PSO-VNS (benchmark: 8; Lillgrund: 18, nine methods at two budgets) and over the 5 sampling-control contrasts (SSA-VNS vs.\ RSD-VNS). All 13 benchmark comparisons significant in the pooled tests stay significant with seeds as the unit. Details: Tables~\ref{S-tab:S-sa-allrun}, \ref{S-tab:S-sa-dep-bench} and~\ref{S-tab:archive-paired}; Section~\ref{S-sec:S-sa-dep}.}
\end{table}

\emph{Coordinate precision.} The stored coordinates are rounded to three (benchmark, IEA37) or two decimals (Horns Rev~1, Lillgrund). Of the 56,540 stored records, the rounded coordinates confirm 50,826 feasibility labels at $10^{-6}$~m and contradict none (Table~\ref{S-tab:S-sa-precision}); rerun with a 17-digit writer, 4,794 of the other 5,714 are confirmed (0 changed). The remaining 920 are SLSQP-based (SciPy's SLSQP is not bit-reproducible across CPU and BLAS builds), among them 220 of the 240 exact-gradient runs; they keep the labels reported by the runs, which the audit neither confirms nor contradicts. Taking every open label as infeasible leaves the order of the eight benchmark methods without GA unchanged (Section~\ref{S-sec:S-sa-feasbounds}). Replayed from the rounded coordinates, 905 of the 51,810 feasible records differ from their stored objective by more than 0.05~pp where the hub-center wake test makes the objective jump, among them 233 of the 21,716 benchmark layouts re-evaluated in Section~\ref{sec:robustness}; ranked by the replayed objective, the eight methods and the site means keep their order (PSO-VNS${}-{}$PSO: $-0.018$~pp, unchanged). All writers of the reproducibility package store 17 significant digits.
% CHECK (manual, 2026-10-04): analysis/rev3_precision_audit.json total (56,540 = 36,190 + 11,640 + 5,500 + 3,210; class i 50,826, class iii 5,714) and rev3_precision_audit_records.csv (Role != duplicate): feasible 51,810, |ReplayDiffPP| > 0.05 for 905; max |ReplayDiffPP| fresh_grid 0.941, mpce_b120k grid 2.698; max |ReplayDiff| mpce_hrfix 0.527 GWh, rev2_lg16b 0.225 GWh; C4 local reruns (analysis/rev3_precision_notes.md): 1,649 records, 1,160 bit-identical, 412 float noise, 11 near + 66 diverged (all SLSQP-based) = 77 undecided, 1,572 confirmed, 0 changed; cloud 4,065 records pending
% CHECK (manual, 2026-10-04; mpce_direction_layouts.csv.gz, J1 = 15-deg objective replayed from the rounded coordinates): |J1 - Objective| > 0.05 pp for 233 of 21,716 layouts (SSA 110, LX-SSA 84, MS-SLSQP 21, SSA-VNS 14, four others 1 each); 8-method feasibility-aware average ranks stored -> replayed: PSOBV 1.816 -> 1.801, PSOC 1.971 -> 1.956, SSABV 3.449 -> 3.485, BVNS 4.397 -> 4.331, SSA 5.382 -> 5.721, SLSQP 5.882 -> 5.750, LXSSA 6.059 -> 6.015, DE 7.044 -> 6.941 (order unchanged); 11 methods: SSA 8.118 -> 8.485 and SLSQP 8.324 -> 8.206 swap; PSO-VNS - PSO -0.0180 -> -0.0179 pp; rev3_sites_reeval.csv: per-budget order of the Horns Rev 1 and Lillgrund method means unchanged under the replay

\subsection{Budget and Initialization}
\label{sec:beyond}
\subsubsection{Budget Scaling}
\label{sec:budget}
% >>>>> begin analysis/mpce_tab_feasbudget.tex
%% generated by mpce_results.py -- do not edit by hand
\begin{table}[!t]
\centering
\caption{Initialization and budget on the six largest benchmark cases (30 seeds).}
\label{tab:feasbudget}
\scriptsize\setlength{\tabcolsep}{2.4pt}
\begin{tabular}{lccccccc}
\toprule
& \multicolumn{2}{c}{Random init.} & \multicolumn{2}{c}{Feasible init.} & \multicolumn{3}{c}{Avg.\ rank} \\
\cmidrule(lr){2-3}\cmidrule(lr){4-5}\cmidrule(lr){6-8}
Method & Loss & Feas. & Loss & Feas. & 6,030 & 30,030 & 120,030 \\
\midrule
PSO-VNS & 4.084 & 100.0 & 3.938 & 100.0 & 1.50 & 1.17 & 2.17 \\
SSA-VNS & 4.678 & 95.0 & 3.955 & 100.0 & 4.33 & 4.17 & 4.17 \\
LX-SSA-VNS & 4.802 & 98.9 & 3.952 & 100.0 & 5.17 & 5.67 & 5.83 \\
RS-VNS & 4.830 & 93.3 & n/r & n/r & 4.67 & 6.33 & 4.67 \\
VNS & 4.950 & 100.0 & 3.810 & 100.0 & 6.50 & 4.83 & 3.17 \\
PSO & 4.181 & 100.0 & 5.549 & 100.0 & 1.83 & 2.50 & 5.83 \\
SSA & 4.828 & 80.0 & 4.947 & 100.0 & 8.33 & 8.67 & 8.83 \\
LX-SSA & 5.379 & 77.8 & 5.048 & 100.0 & 8.67 & 9.83 & 10.00 \\
DE & -- & 1.1 & 5.549 & 100.0 & 10.00 & 5.17 & 3.33 \\
MS-SLSQP & 4.522 & 99.4 & 4.294 & 100.0 & 4.00 & 6.67 & 7.00 \\
\bottomrule
\end{tabular}
\par\vspace{2pt}\parbox{\columnwidth}{\scriptsize Loss: mean wake loss (\%) of the feasible runs; Feas.: feasible runs (\%); both at 6,030 evaluations with random and feasibility-preserving initialization. Avg.\ rank at 6,030, 30,030 and 120,030 evaluations, random initialization. n/r: not run. GA was not run in this study.}
\end{table}
% <<<<< end analysis/mpce_tab_feasbudget.tex

The budget study covers the six largest benchmark cases, the nine methods except GA and the ablation variants LX-SSA-VNS and RS-VNS (Table~\ref{tab:feasbudget}; Fig.~\ref{fig:budget}). From random starts:
\begin{itemize}
\item PSO-VNS keeps the best average rank at all three budgets (6,030, 30,030 and 120,030 evaluations: 1.50, 1.17 and 2.17) and ranks first in 4, 5 and 3 of the six cases (rank differences not tested); PSO's mean wake loss minus that of PSO-VNS is +0.10, +0.18 and +0.27~pp (Table~\ref{S-tab:budget-cases}).
\item The methods differ mainly in how much they gain from a larger budget: from 6,030 to 120,030 evaluations, PSO falls back from an average rank of 1.83 to 5.83, behind SSA-VNS (4.17), and VNS moves up from 6.50 to 3.17; the mean wake loss of VNS falls from 4.95\% to 3.17\%, that of PSO only from 4.18\% to 3.42\% (PSO-VNS: 4.08\% to 3.14\%).
\item At 120,030 evaluations, VNS, SSA-VNS, RS-VNS, LX-SSA-VNS and DE come within 0.17~pp of the mean wake loss of PSO-VNS.
\end{itemize}
% CHECK-FINAL [C26] [C29] [C34]: PSO-VNS best average rank at all three budgets; first-place counts generated per budget
% CHECK-FINAL [C37]: at 120,030 SSA-VNS ranks ahead of PSO; VNS, SSA-VNS, RS-VNS, LX-SSA-VNS and DE within 0.2 pp (\NBudCloseGapOneTwentyK) of PSO-VNS. PSO and VNS ranks are generated macros (not tested by C37).
% verified (manual): budget losses from mpce_numbers.tex (\NBudLoss...): VNS - PSO-VNS mean wake loss 0.87 / 0.40 / 0.03 pp at 6,030 / 30,030 / 120,030 (decreasing); drop 6,030 -> 120,030: VNS 1.78, PSO 0.76, PSO-VNS 0.94 pp; DE feasible share at 6,030 from Table feasbudget (1.1 %); DE plotted only at 30,030 and 120,030 (at least half of the runs feasible)

\begin{figure}[!htbp]
\centering
\pipefig{0.64\textwidth}{figures_mpce/budget_scaling.pdf}{6cm}
\caption{Budget scaling: mean wake loss of the six largest benchmark cases and mean AEP loss of the Horns Rev~1 block (feasible runs; methods with at least half of their runs feasible) versus the budget (random initialization; 10 seeds for Horns Rev~1 at 120,030 evaluations).}
\label{fig:budget}
\end{figure}

On the Horns Rev~1 block (last panel of Fig.~\ref{fig:budget}), the mean AEP loss of PSO-VNS also falls with the budget (7.07\%, 6.52\% and 6.40\%), and at 120,030 evaluations DE has the lowest.
% CHECK-FINAL [C25]: Horns Rev 16: 30,030 and 120,030 mean AEP loss of PSO-VNS below 6,030; CHECK-DIR [F29]: DE highest mean AEP (lowest loss) at 120,030

\subsubsection{Feasibility-Preserving Initialization}
\label{sec:feasinit}
With feasibility-preserving initialization (Table~\ref{tab:feasbudget}; per case: Table~\ref{S-tab:feasinit-cases}), the share of feasible runs rises to 100.0\% for SSA-VNS (from 95.0\%), SSA (80.0\%), LX-SSA (77.8\%) and DE (1.1\%). VNS alone is then best in mean wake loss and average rank, and PSO-VNS ranks second (rank 2.83; loss 3.94\%, 4.08\% with random starts). PSO and DE never improve on the best initial layout (210 runs each; Table~\ref{S-tab:D-feasstart}): their moves change all or most coordinates at once, so in our runs no trial is both feasible and better and the global best never moves (Lemma~\ref{lem:feas-improve}, Proposition~\ref{prop:feasible-start}); VNS, which moves one coordinate at a time, is not affected. The advantage of PSO-VNS thus depends on the initialization: from random starts, PSO may help mainly by reaching feasibility.
% CHECK-FINAL [C28] [C38]: feasible starts: VNS alone lowest loss and best rank; PSO-VNS second in rank; PSO and DE never move from the best initial layout
% CHECK-DIAG [D13]: PSO and DE never improve on the best initial layout (instrumented and all stored runs)
% CHECK-DIAG [D14]: all PSO candidates except the G re-evaluations and all \NDFsCandDE{} DE trials infeasible, none accepted, no pbest update (first-move median min spacing \NDFsFirstMinSpPSO{} m: printed in the diagnostics paragraph below)

What follows from the algorithms can be separated from what is only observed. Let $L(\mathbf X)$ be the wake loss, the first term of~\eqref{eq:penalty}, so that $F_p=L+\mathrm{Pen}$ with $\mathrm{Pen}\ge0$.
\begin{lemma}[Improvement needs a feasible, better trial]
\label{lem:feas-improve}
If the global best $\mathbf G$ is feasible, a trial $\mathbf X$ replaces it only if $L(\mathbf X)<L(\mathbf G)$ and every constraint violation $g$ of $\mathbf X$ satisfies $(1+\mu_{\rm p} g)^2<L(\mathbf G)$; in the benchmark cases, $g<4.6\times10^{-8}$ (m or m$^2$), far below the reporting tolerance of $10^{-6}$~m.
\end{lemma}
\begin{proof}
Replacement needs $F_p(\mathbf X)<F_p(\mathbf G)=L(\mathbf G)$, and $F_p(\mathbf X)$ is the sum of $L(\mathbf X)\ge0$ and the violation terms $(1+\mu_{\rm p} g)^2>0$. In the benchmark cases, $L(\mathbf G)\le NE_{\rm ideal}+\varepsilon$ with $N\le15$, $E_{\rm ideal}\le14{,}045.74$ ($\varepsilon$: the small negative part of the power curve above cut-in), so $\sqrt{L(\mathbf G)}<459.1$ and $g<(\sqrt{L(\mathbf G)}-1)/\mu_{\rm p}<4.6\times10^{-8}$ for $\mu_{\rm p}=10^{10}$.
\end{proof}
% verified: (sqrt(15 x 14045.74) - 1)/1e10 = 4.58e-8 (as optA/supp_theory.tex, Lemma S-feas-improve; the small negative part of the power curve above cut-in, epsilon, does not change the bound)
\begin{proposition}[PSO from feasible starts]
\label{prop:feasible-start}
Let all initial layouts be feasible, let the velocities start at zero with $\mathbf P^i=\mathbf X^i$, and let $i^*$ be the particle with the best initial layout, so that $\mathbf X^{i^*}=\mathbf P^{i^*}=\mathbf G$.
(a)~As long as $\mathbf G$ does not change, particle $i^*$ keeps $\mathbf V^{i^*}=\mathbf 0$ and $\mathbf X^{i^*}=\mathbf G$: it re-evaluates $\mathbf G$ once per iteration and evaluates no other point.
(b)~The first move of every other particle is $\mathbf X^i\leftarrow\mathbf X^i+c_2\boldsymbol\rho_2\circ(\mathbf G-\mathbf X^i)$, clipped to the box: turbine $k$ of layout $i$ moves a random fraction in $[0,c_2]=[0,1.50]$ of the way, per coordinate, toward turbine $k$ of $\mathbf G$, although the numbering of the turbines of two independently generated layouts is arbitrary.
\end{proposition}
\begin{proof}
(a)~By induction: with $\mathbf V^{i^*}=\mathbf 0$ and $\mathbf P^{i^*}=\mathbf X^{i^*}=\mathbf G$, all three terms of~\eqref{eq:pso-v} vanish, the particle stays at $\mathbf G$, and $F_p(\mathbf G)$ is no strict improvement, so neither $\mathbf P^{i^*}$ nor $\mathbf G$ changes. (b)~At the first move, $\mathbf V^i=\mathbf 0$ and $\mathbf P^i=\mathbf X^i$, so only the last term of~\eqref{eq:pso-v} remains.
\end{proof}
A DE trial for population member $i$ takes coordinate $j$ from the mutant $\mathbf X^{r_1}+F(\mathbf X^{r_2}-\mathbf X^{r_3})$, with distinct random members $r_1,r_2,r_3\ne i$, if $j=j_{\rm rand}$ (a random coordinate) or $U_j<\mathrm{CR}$ ($U_j\sim U(0,1)$), i.e., $1+\mathrm{Bin}(n-1,\mathrm{CR})$ coordinates, on average 27.1 of $n=30$ for $N=15$, again regardless of the turbine numbering. By Lemma~\ref{lem:feas-improve}, $\mathbf G$ moves only if such a combined layout is feasible and better; that this never happened in our runs is an observation, not a consequence of the proposition.
% CHECK-THEORY [T11]: DE takes most coordinates from the mutant (mean 1+CR(n-1) = 27.1 of 30 for N = 15, CR = 0.9; PSO moves all). Lemma lem:feas-improve = Lemma S-lem:S-feas-improve and Proposition prop:feasible-start = Proposition S-prop:S-frozen (a), (b) (optA/supp_theory.tex, S-th-feas); the stall is observed (D13/D14), not implied.
% Page budget: kept in the supplementary material: \SuppTable{tab:D-feasstart}

In Table~\ref{S-tab:D-feasstart}, which logs every candidate, the only feasible PSO candidates (3.33\%) are the re-evaluations of $\mathbf G$ of part~(a), no personal best is updated, and no DE trial is feasible. The numbering of part~(b) is not the whole explanation: after relabeling the turbines of each start to their nearest turbines of $\mathbf G$, the first PSO move is feasible in only 0.1\% of the draws, a convex combination of two layouts in 16.6\%; independent weights per coordinate, with overshoot, place the turbines of two distinct dense layouts (median RMS distance 291~m) too close together (median minimum spacing after the first move 95~m, limit 308~m) or outside the site.
% CHECK-DIAG [D16]: the feasible initial layouts are distinct (matched RMS distance to G > 50 m)
% CHECK-DIAG [D14]: every feasible PSO candidate is the G-particle (zero velocity) re-evaluating G; no personal best updated; DE no feasible trial, none accepted; first-move median min spacing \NDFsFirstMinSpPSO{} m (95 m)
% CHECK-DIAG [D15]: after optimal relabelling the first PSO move is feasible in < 1 % of draws; only a common weight (convex combination) more often (\NDFsMixScalar %)
% CHECK-DIAG [D16]: feasible initial layouts distinct (matched RMS distance to G > 50 m; median \NDFsInitRms m)

\subsection{Sensitivity to the Model Assumptions}
\label{sec:robustness}
% tab:robust-final is placed in the main text; the full re-evaluation stays in the supplement (Table S-robust). tab:F-bench and tab:F-equiv are
% captured in 06_type1.tex (\CaptureSuppTables{analysis/mpce_supp_direction.tex}).
The benchmark (Type~I) fixes five modeling choices. We tested each of them with the nine methods:
\begin{enumerate}
\item \emph{power curve} and
\item \emph{wake model}: every feasible final layout re-evaluated, not re-optimized (Table~\ref{tab:robust-final}; all methods: Table~\ref{S-tab:robust}), which asks whether the delivered layouts keep their order under another model;
\item \emph{minimum spacing}: methods re-optimized at $5D$ and $6D$ on the split cases;
\item \emph{boundary rule}: five methods re-optimized under two other constraint handlings on six cases;
\item \emph{direction resolution}: final layouts re-evaluated with finer wind roses, and the fine rose optimized directly on eight cases (Table~\ref{tab:direction}).
\end{enumerate}
Re-evaluations use the stored rounded coordinates (Section~\ref{sec:archive-reliability}).
% >>>>> begin analysis/rev5_ga_tab_robust_final.tex
%% generated by rev5_ga_reeval.py (pipeline of mpce_results.py section 8, GA added to the pool) -- do not edit by hand
\begin{table}[!t]
\centering
\caption{Robustness to the benchmark model (Type~I, nine methods): final layouts re-evaluated, not re-optimized.}
\label{tab:robust-final}
\scriptsize\setlength{\tabcolsep}{2.4pt}
\begin{tabular}{lcccccc}
\toprule
& \multicolumn{2}{c}{$\Delta$ (\%)} & & & & \\
\cmidrule(lr){2-3}
Model & DS I & DS II & Rank & Best & $\bar\tau$ & Same \\
\midrule
Benchmark & +0.0 & +0.0 & 2.18 & PSO-VNS & -- & -- \\
Cubic curve & $-17.7$ & $-36.9$ & 2.24 & PSO-VNS & 0.94 & 90 \\
Cubic, 25~m/s cut-out & $-21.5$ & $-37.0$ & 2.22 & PSO-VNS & 0.94 & 90 \\
Gaussian wake & $-0.3$ & $-0.4$ & 2.31 & PSO & 0.66 & 49 \\
\bottomrule
\end{tabular}
\par\vspace{2pt}\parbox{\columnwidth}{\scriptsize All feasible final layouts of the 68 cases. $\Delta$: mean relative change of the objective (\%); Rank: average rank of PSO-VNS; Best: best-ranked method; $\bar\tau$: mean Kendall correlation between the benchmark and the alternative ordering within a case; Same: cases (\%) with unchanged best method. Ranks of all methods: supplementary material.}
\end{table}
% <<<<< end analysis/rev5_ga_tab_robust_final.tex

\subsubsection{Power Curve}
\label{sec:powercurve}
A cubic power curve between cut-in and rated speed~\cite{Carrillo2013} (also with a 25~m/s cut-out) lowers the expected power by 17.7\% (Data Set~I) and 36.9\% (Data Set~II), so our energy values are benchmark values, but it leaves the ranking intact ($\bar\tau=0.94$; same best method in 90\% of the cases; PSO-VNS first, PSO second and GA third, also with the cut-out). The $\pm0.05$~pp equivalence verdicts (Sections~\ref{sec:psovns-pso} and~\ref{sec:ablation}) were evaluated under the linear benchmark curve only.
% CHECK-FINAL [C21]: robustness.Cubic/CubicCutout: tau >= 0.9, PSO-VNS rank position 1

\subsubsection{Wake Model}
\label{sec:gaussian}
The uncalibrated Gaussian wake of~\cite{Bastankhah2014} ($K_{\rm G}=0.04$; Section~\ref{S-sec:S-gauss}) differs from the Jensen cone ($K=0.075$) in its smooth lateral profile and expansion rate and gives no deficit to turbines upstream of the wake source. It changes the wake loss of the same final layouts by 0.38~pp on average (median 0.28~pp), 7.6 times the equivalence margin (Section~\ref{sec:setup-stats}), and by more than the margin for 86\% of the 17{,}626 layouts, whereas the per-case difference PSO-VNS${}-{}$PSO changes by only 0.031~pp on average. The ordering is less stable ($\bar\tau=0.66$): PSO is narrowly first (average rank 2.24; PSO-VNS 2.31, GA 3.09), but not with 1$^\circ$ bins (Table~\ref{tab:direction}).
% CHECK-EXTRA [X30]: Jensen vs Gaussian wake loss of the same layouts: mean \NXModelShiftMean pp (median \NXModelShiftMedian; \NXModelShiftRatio x margin; > margin for \NXModelShiftAboveMarginPct % of \NXModelShiftNLayouts layouts); per-case PSO-VNS - PSO changes by \NXModelShiftPairMean pp on average (below the margin)
% CHECK-FINAL [C22]: robustness.Gauss (15 deg): tau < 0.8, PSOC first, PSOBV second (nine-method pool with GA: 2.24 vs 2.31, rev5_ga_numbers_robust.tex); CHECK-DIR [F09]: PSO-VNS first under the Gaussian wake at 1 deg

\subsubsection{Minimum Spacing}
\label{sec:spacing}
The optimized layouts often lie on the $4D$ constraint: 49\% of the feasible final layouts also satisfy $5D$ and 36\% $6D$ (4.1\% and 0.1\% for $N=11$--15, where the fixed site leaves less room), so they do not transfer to larger spacings without re-optimization (site capacity: Table~\ref{S-tab:capacity}).
% CHECK-FINAL [C23]: spacing_share.n11_15.pct_5d < 10 ("often lie on the 4D constraint")
% verified (manual): shares 49 % / 36 % (all N) vs 4.1 % / 0.1 % (N = 11-15) from mpce_numbers.tex (\NSp...); constructive capacity at 4D/5D/6D: Table S-capacity

The eight methods without GA and RSD-VNS were re-optimized at $5D$ and $6D$ on the 12 split cases, of which the two with $r=500$~m, $N=10$ admit no layout at these spacings, so the same 10 cases are compared at every spacing (Section~\ref{S-sec:S-xp-spacing}; Table~\ref{S-tab:S-xp-spacing}):
\begin{itemize}
\item PSO-VNS keeps the best average rank (1.40, 1.60 and 1.90 at $4D$, $5D$ and $6D$); its mean advantage over PSO is larger at $5D$ and $6D$ ($-0.094$ and $-0.068$~pp, the latter over the 8 cases in which PSO has at least 15 feasible runs) than at $4D$ ($-0.035$~pp), with 90\% intervals below zero at the seed and the case level at $5D$ and $6D$.
\item MS-SLSQP moves up with the spacing (average rank 5.50, 4.30 and 3.00; second at $6D$).
\item The SSA phase gives no gain over disc sampling at $5D$ (SSA-VNS${}-{}$RSD-VNS: $+0.043$~pp, case-level 90\% CI [$-$0.011, 0.103]), but at $6D$ SSA-VNS has the lower mean wake loss in all 8 cases in which both qualify ($-0.136$~pp, [$-$0.215, $-$0.073]; Wilcoxon $p=0.0078$): the no-gain finding of Section~\ref{sec:ablation} holds at $4D$ and $5D$, not at $6D$.
\item On the all-run endpoint (Table~\ref{S-tab:S-allrun-spacing}), PSO-VNS scores 0.593 and 0.613 against PSO at $5D$ and $6D$ (95\% case-bootstrap intervals [0.497, 0.693] and [0.548, 0.670]), but 0.580 [0.432, 0.733] against MS-SLSQP at $6D$: no advantage.
\end{itemize}
With 8 to 10 cases, these are estimates for the split cases, not general effects.

\subsubsection{Boundary Rule and Constraint Handling}
\label{sec:boundary}
All methods clip the coordinates to the bounding square and leave the circle to the penalty (Section~\ref{sec:constraints}). The rule matters: an SSA that projects turbines radially onto the circle instead attains higher mean objectives than the original SSA in all six representative cases ($p\le0.0036$).
% CHECK-FINAL [C41]: boundary_rule n_higher == n_cases == 6 and max_p < 0.01
PSO-VNS, PSO, GA, SSA-VNS and DE were therefore rerun on six cases (seeds 31--60, 6,030 evaluations; Table~\ref{S-tab:S-sa-constraint}) under Deb's feasibility rules on the dimensionless violations $g^{\rm b}_i/r^2$ and $g^{\rm s}_{ij}/\ell_{\min}$ (reference runs under the penalty reproduce the stored runs bit for bit):
\begin{itemize}
\item \emph{With radial projection}, the conditional wake loss of every method fell by 0.46--0.80~pp and the feasibility of DE rose from 33.3\% to 80.0\%; the run-level order PSO-VNS, PSO, GA, SSA-VNS, DE and the leading case-level method (PSO-VNS) persisted; PSO-VNS kept its all-run advantage over GA, SSA-VNS and DE (0.683, 0.850, 0.844), but not over PSO (0.561; conditional $\overline{\Delta L}=-0.095$~pp [$-0.162$, $-0.037$]).
\item \emph{With box clipping}, the feasibility of PSO fell from 98.9\% to 73.3\%, the conditional loss of PSO-VNS rose by 0.23~pp, its all-run advantage over GA (0.517) and SSA-VNS (0.583) was not significant, and GA had the best case-level rank (1.83 vs.\ 2.17; run level: PSO-VNS 2.19, GA 2.28).
\end{itemize}
With the 30 seeds as the unit of replication, every all-run verdict is unchanged (Table~\ref{S-tab:S-sa-dep-studies}). By the survival rule stated beforehand, the conclusions fail with box clipping and, for the all-run advantage of PSO-VNS over PSO, with projection: the comparison is one of search strategies under a common penalty and repair, which is not neutral across methods.
% CHECK (manual, 2026-10-04): constraint study (analysis/rev3_constraint.json, rev3_constraint_tables.tex, tab:S-sa-constraint): variant (i) reproduces 60/60 stored runs bit for bit; (iii) losses (i)->(iii) PSO-VNS 3.126->2.332, PSO 3.237->2.460, GA 3.307->2.550, SSA-VNS 3.502->3.044 (0.46-0.80 pp); DE feas 33.3->80.0; run-level order unchanged; scores 0.683/0.850/0.844 (p_H < 0.001), PSO 0.561 (p_H 1.000), dL -0.095 [-0.162, -0.037]; (ii) PSO feas 98.9 -> 73.3, PSO-VNS loss 3.126 -> 3.353 (+0.23), scores GA 0.517 (p_H 1.000), SSA-VNS 0.583 (p_H 0.395); case-level rank GA 1.83 vs PSO-VNS 2.17, run level 2.19 vs 2.28; survival rule (manifest): fails for (ii), fails only on PSO for (iii)

\subsubsection{Direction Resolution}
\label{sec:direction}
Because a pair of turbines can be placed with no bin center in its narrow wake interval (Section~\ref{sec:wind}), we split each bin into 3 (5$^\circ$) or 15 (1$^\circ$) sub-bins with its Weibull parameters and an equal share of its frequency (wake-free power unchanged) and re-evaluated all 23{,}728 feasible final layouts of the benchmark methods (Tables~\ref{tab:direction} and~\ref{S-tab:F-bench}).
\begin{table}[!htb]
\centering
\caption{Direction resolution on Type~I: average rank of the three leading methods when the final layouts are re-evaluated with finer wind roses (nine methods, 68 cases) and when the 1$^\circ$ objective is optimized directly (five methods, eight cases).}
\label{tab:direction}
\footnotesize\setlength{\tabcolsep}{4pt}
\begin{tabular}{llcccl}
\toprule
Wind rose & Wake model & PSO-VNS & PSO & GA & Best \\
\midrule
15$^\circ$ (benchmark) & Jensen & 2.18 & 2.40 & 2.74 & PSO-VNS \\
1$^\circ$, re-evaluated & Jensen & 2.26 & 2.57 & 4.01 & PSO-VNS \\
15$^\circ$ (benchmark) & Gaussian & 2.31 & 2.24 & 3.09 & PSO \\
1$^\circ$, re-evaluated & Gaussian & 2.49 & 2.56 & 3.96 & PSO-VNS \\
1$^\circ$, optimized (8 cases) & Jensen & 1.12 & 2.62 & 3.38 & PSO-VNS \\
\bottomrule
\end{tabular}
\par\vspace{2pt}\parbox{\columnwidth}{\scriptsize Re-evaluated rows: feasibility-aware average rank among the nine methods (all ranks: Table~\ref{S-tab:F-bench}). Optimized row: PSO-VNS, PSO, GA, MS-SLSQP and RSD-VNS optimized with 1$^\circ$ bins on eight cases, seeds 31--60, 6,030 evaluations (Table~\ref{S-tab:S-sa-fine}).}
\end{table}
\begin{itemize}
\item The 1$^\circ$ rose raises the mean wake loss of every method, by 0.40~pp (DE) to 1.09~pp (GA; PSO-VNS: 1.322\% to 2.335\%), mostly already with 5$^\circ$ sub-bins: optimized layouts exploit the gaps between the coarse bins, those of GA, PSO and the VNS-based methods most.
\item PSO-VNS still ranks first, also under the Gaussian wake at 1$^\circ$; PSO is the only method not significantly different. Beyond the top two the order changes ($\bar\tau=0.50$; same best method in 31\% of the cases): GA falls to third with the largest loss rise, and MS-SLSQP moves from seventh to fourth; its gap to PSO-VNS shrinks from 0.435 to 0.135~pp but remains significant.
\item PSO-VNS has a lower case-mean wake loss than GA at 5$^\circ$ and 1$^\circ$ ($-0.098$ and $-0.105$~pp, Holm $p\le3.1\times10^{-7}$), whereas at 15$^\circ$ the difference ($-0.026$~pp) is not significant.
\item For PSO-VNS vs.\ PSO (details: Table~\ref{S-tab:F-equiv}), the case-mean difference becomes $-0.043$~pp (90\% CI [$-0.062$, $-0.025$]; $p=2.5\times10^{-4}$; the same with 5$^\circ$ sub-bins): a small but significant advantage, not equivalence at $\pm0.05$~pp (minimal margin 0.062~pp). On the 44 non-trivial cases, equivalence is shown at neither resolution, and the Data Set~II $N\ge10$ advantage remains (10 of 10 cases).
\item The ablation conclusions hold, but the small SSA gain over RS-VNS is not significant with 1$^\circ$ bins (Table~\ref{S-tab:F-abl}).
\end{itemize}
The order beyond the top two and the equivalence verdict are thus properties of the discretized benchmark; the first place of PSO-VNS is not.
% CHECK-DIR [F03]: loss rise of every method, PSO-VNS 1.322 -> 2.335; [F04]: most of the change already with 5-deg sub-bins (> 70 % for PSO-VNS); [F05]: VNS-based methods and PSO rise most, DE/MS-SLSQP least ("VNS-based methods most"); [F06]: PSO-VNS first at 5 and 1 deg, PSO the only one not significantly different (Holm); [F09]: Gaussian: PSO first at 15 deg, PSO-VNS first at 1 deg; [F07]: MS-SLSQP 7th -> 4th with GA in the pool (6th -> 3rd without GA), tau 0.52, same best 37 %; [F08]: MS-SLSQP gap shrinks from 0.435 to 0.135 pp, still significant (Holm)
% CHECK-DIR [F10]: PSO-VNS - PSO -0.043, CI [-0.062, -0.025], not equivalent, p < 0.05; [F11]: same with 5-deg sub-bins; [F12]: per-case sign changes in 26 of 60 cases; [F13]: non-trivial cases: equivalence at neither resolution; [F14]: DS II N >= 10 10/10; [F16]: component-analysis conclusions hold (SSA-VNS and LX-SSA-VNS worse than RSD-VNS, PSO-VNS better than RSD-VNS, RS-VNS and VNS, PSO-VNS best of nine); [F17]: SSA-VNS vs RS-VNS not significant at 1 deg
% Page budget: kept in the supplementary material: \SuppTable{tab:F-bench}

\emph{Direct optimization with 1$^\circ$ bins.} PSO-VNS, PSO, GA, MS-SLSQP and RSD-VNS also optimized the 1$^\circ$ objective directly on eight cases (both data sets, $N=6$--15; seeds 31--60; 6,030 evaluations), compared with a matched control arm optimized with 15$^\circ$ bins and evaluated with 1$^\circ$ bins (Table~\ref{S-tab:S-sa-fine}).
\begin{itemize}
\item Optimizing with 1$^\circ$ bins lowers the 1$^\circ$ wake loss on the same seeds by 0.84--1.12~pp for the metaheuristics and by 0.32~pp for MS-SLSQP, which drops from second to last.
\item PSO-VNS ranks first in both arms (direct: 1.12; PSO 2.62, GA 3.38, RSD-VNS 3.62, MS-SLSQP 4.25; control: 2.25).
\item In the direct arm it has a lower wake loss than PSO in all eight cases ($-0.194$~pp; 90\% CI [$-0.234$, $-0.152$] at the seed and [$-0.272$, $-0.114$] at the case level), beyond the margin and about twice its advantage in the control arm ($-0.086$~pp; difference-in-differences on the 238 of 240 case--seed pairs where all four runs are feasible: $-0.115$~pp, exact sign-flip $p=3.9\times10^{-4}$; Table~\ref{S-tab:S-sa-dep-studies}).
\item Its all-run scores against the other four methods in the direct arm are 0.70--0.87, all significant with the 30 seeds as the unit of replication ($p_{\rm Holm}\le3.2\times10^{-6}$).
\end{itemize}
The leader thus does not change when the search uses 1$^\circ$ bins.
% CHECK (manual, 2026-10-04): direct 1-deg study (analysis/rev3_fine.json, rev3_fine_tables.tex, tab:S-sa-fine): direct avg ranks PSOBV 1.125, PSOC 2.625, GA 3.375, RSDVNS 3.625, SLSQP 4.25; PSO-VNS - PSO -0.194 pp, seed 90% [-0.234, -0.152], case 90% [-0.272, -0.114], Wilcoxon p 0.0078 (8 of 8 cases); all-run scores 0.696 (RSD-VNS) to 0.871 (MS-SLSQP), sign-test p <= 1.2e-9; control arm (seeds 31-60, 15-deg optimized, 1-deg re-evaluated): ranks PSOBV 2.25, SLSQP 2.625, PSOC 2.875, GA 3.125, RSDVNS 4.125, PSO-VNS - PSO -0.086 pp; 1-deg loss lowered by direct optimization on the same seeds: PSO-VNS 1.12, PSO 1.01, RSD-VNS 1.09, GA 0.84, MS-SLSQP 0.32 pp
% Over the page budget (+0.5 pages in the private build), kept in the supplement: \SuppTable{tab:F-equiv}
